\documentclass[main.tex]{subfiles}
\begin{document}

% ╔══════════════════════════════════════════════════════════════════════════════╗
% ║  showcase.tex — Every feature demonstrated                                   ║
% ║  Inspect this to see what every box, list, table, and command looks like.    ║
% ╚══════════════════════════════════════════════════════════════════════════════╝

\chapter{Feature Showcase — All Boxes}

This chapter demonstrates every single environment and command from the preamble.
Use it to inspect the output and as a copy-paste reference for your own notes.

% ═══════════════════════════════════════════════════════════════════════════════
\section{All Eight Box Types}
% ═══════════════════════════════════════════════════════════════════════════════

% ─── 1. Definition Box ───────────────────────────────────────────────────────
\begin{definitionbox}[Network Protocol]
A \term{network protocol} is a set of rules and conventions that govern how
data is transmitted between devices on a network. Protocols define:
\begin{itemize}
    \item The format and structure of messages
    \item The sequence in which messages are sent and received
    \item How errors are detected and corrected
    \item How connections are established and terminated
\end{itemize}
\end{definitionbox}

\vspace{6pt}

% ─── 2. Formula Box ──────────────────────────────────────────────────────────
\begin{formulabox}[Shannon's Channel Capacity]
The maximum data rate of a noisy channel (Shannon--Hartley theorem):
\[
    C = B \log_2\!\left(1 + \frac{S}{N}\right)
\]
where $C$ is channel capacity in bits/s, $B$ is bandwidth in Hz,
and $S/N$ is the signal-to-noise ratio (linear, not dB).

\Note For $S/N = 0$ (pure noise), $C = 0$ regardless of bandwidth.
\end{formulabox}

\vspace{6pt}

% ─── 3. Info Box ─────────────────────────────────────────────────────────────
\begin{infobox}[Analogy: Postal System]
Think of the OSI model layers like the postal system:
\begin{itemize}
    \item The \term{Application Layer} is you writing the letter.
    \item The \term{Transport Layer} is the envelope with the address.
    \item The \term{Network Layer} is the postal routing system.
    \item The \term{Physical Layer} is the actual truck carrying the mail.
\end{itemize}
Each layer only talks to the layer directly above and below it.
\end{infobox}

\vspace{6pt}

% ─── 4. Derivation Box ───────────────────────────────────────────────────────
\begin{derivationbox}[Deriving the Propagation Delay Formula]
Start with the definition of propagation speed $v$ along a medium:
\[
    v = \frac{d}{t_p}
\]
where $d$ is distance and $t_p$ is propagation delay. Rearranging:
\[
    t_p = \frac{d}{v}
\]
For copper wire, $v \approx 2 \times 10^8\ \unit{m/s}$ (roughly $\tfrac{2}{3}c$).

For a 1\,km cable:
\[
    t_p = \frac{1000\ \unit{m}}{2 \times 10^8\ \unit{m/s}} = 5\ \unit{\mu s}
\]
\end{derivationbox}

\vspace{6pt}

% ─── 5. Example / Q&A Box (auto-numbered) ────────────────────────────────────
\begin{questionanswerbox}
\textbf{Problem:} A channel has a bandwidth of $B = 4\ \unit{kHz}$ and SNR $= 31$.
Find the maximum channel capacity using Shannon's theorem.

\medskip
\textbf{Solution:}
\[
    C = B \log_2(1 + \text{SNR})
      = 4000 \times \log_2(32)
      = 4000 \times 5
      = 20{,}000\ \unit{bits/s}
      = 20\ \unit{kbps}
\]
\end{questionanswerbox}

\vspace{4pt}

\begin{questionanswerbox}
\textbf{Problem:} Compute the Nyquist capacity for a noiseless channel with
$B = 3\ \unit{kHz}$ and $M = 8$ signal levels.

\medskip
\textbf{Solution:}
\[
    C_{\text{Nyquist}} = 2B\log_2 M = 2 \times 3000 \times \log_2 8 = 18{,}000\ \unit{bps}
\]
\end{questionanswerbox}

\vspace{6pt}

% ─── 6. Alert Box ────────────────────────────────────────────────────────────
\begin{alertbox}[Common Mistake: dB vs Linear SNR]
Shannon's formula requires SNR as a \textbf{linear ratio}, not in decibels.

\medskip
If given $\text{SNR}_{\text{dB}} = 30\ \unit{dB}$, convert first:
\[
    \text{SNR}_{\text{linear}} = 10^{30/10} = 10^3 = 1000
\]
\Imp Never plug the dB value directly into the $\log_2$ formula.
\end{alertbox}

\vspace{6pt}

% ─── 7. Summary Box ──────────────────────────────────────────────────────────
\begin{summarybox}[Section Recap: Channel Capacity]
\begin{itemize}
    \item \textbf{Nyquist (noiseless):} $C = 2B\log_2 M$ — depends on signal levels
    \item \textbf{Shannon (noisy):} $C = B\log_2(1+S/N)$ — hard theoretical limit
    \item Increasing bandwidth $B$ always helps, but diminishing returns apply
    \item SNR must be linear (not dB) in Shannon's formula
\end{itemize}
\end{summarybox}

\vspace{6pt}

% ─── 8. Theorem Box ──────────────────────────────────────────────────────────
\begin{theorembox}[Nyquist's Sampling Theorem]
If a continuous signal contains no frequency components higher than $f_{\max}$,
then it can be completely reconstructed from samples taken at a rate of at least
$2f_{\max}$ samples per second (the \emph{Nyquist rate}).
Mathematically:
\[
    f_s \geq 2 f_{\max}
\]
Sampling below this rate causes \emph{aliasing} — high-frequency components fold
back into the lower-frequency range and the signal is unrecoverable.
\end{theorembox}

% ═══════════════════════════════════════════════════════════════════════════════
\section{All Six Heading Levels — Section — Purple}
% ═══════════════════════════════════════════════════════════════════════════════

This section demonstrates all heading levels so you can see every color.

\subsection{Subsection — Neon magenta}

Content under a subsection.

\subsubsection{Subsubsection — Neon Blue}

Content under a subsubsection.

\paragraph{Paragraph — Warm Orange}

Content under a paragraph heading.

\subparagraph{Subparagraph — Mint Teal}

Content under a subparagraph heading. These are the six heading levels available.

% ═══════════════════════════════════════════════════════════════════════════════
\section{Lists — All Nesting Levels}
% ═══════════════════════════════════════════════════════════════════════════════

\subsection{Itemize — 4 Levels}

\begin{itemize}
    \item Level 1 — purple triangle ($\triangleright$)
    \begin{itemize}
        \item Level 2 — blue arrow ($\rightarrow$)
        \begin{itemize}
            \item Level 3 — orange dash (\textendash)
            \begin{itemize}
                \item Level 4 — teal circle ($\circ$)
            \end{itemize}
        \end{itemize}
    \end{itemize}
    \item Another Level 1 item
    \item Yet another Level 1 item
\end{itemize}

\subsection{Enumerate — 3 Levels}

\begin{enumerate}
    \item First item (purple arabic numerals)
    \begin{enumerate}
        \item Sub-item (blue letters)
        \begin{enumerate}
            \item Sub-sub-item (orange roman numerals)
            \item Second roman numeral
        \end{enumerate}
        \item Another sub-item
    \end{enumerate}
    \item Second item
    \item Third item
\end{enumerate}

\subsection{Mixed List Inside a Box}

\begin{infobox}[Lists Work Inside Boxes Too]
You can freely mix list types inside any box:
\begin{itemize}
    \item OSI Model layers (top to bottom):
    \begin{enumerate}
        \item Application
        \item Presentation
        \item Session
        \item Transport
        \item Network
        \item Data Link
        \item Physical
    \end{enumerate}
    \item Each layer adds its own header/trailer when encapsulating data.
\end{itemize}
\end{infobox}

% ═══════════════════════════════════════════════════════════════════════════════
\section{Longtables}
% ═══════════════════════════════════════════════════════════════════════════════

% Local table spacing example:
% Wrap a table in braces to change row height and cell padding for that table only.
%
% {
% \renewcommand{\arraystretch}{1.6}
% \setlength{\tabcolsep}{8pt}
%
% \begin{longtable}{|L{3cm}|L{5cm}|}
%     \longtableheader{2}{Local Spacing Example}{%
%       \tblcolhdcell{Column A} & \tblcolhdcell{Column B}}
%     Row title & Row content \\ \tblhline
%     Another row & More content \\
%     \tblbottomrule
% \end{longtable}
% }

\begin{longtable}{|L{3.5cm}|C{2.5cm}|C{2cm}|L{4.5cm}|}
    \longtableheader{4}{Common Network Protocols}{%
      \tblcolhdcell{Protocol} & \tblcolhdcell{Layer} & \tblcolhdcell{Port} & \tblcolhdcell{Description}}
    HTTP  & Application & 80  & Hypertext Transfer Protocol \\ \tblhline
\evenrow
    HTTPS & Application & 443 & HTTP over TLS/SSL \\ \tblhline
    FTP   & Application & 21  & File Transfer Protocol \\ \tblhline
\evenrow
    SSH   & Application & 22  & Secure Shell \\ \tblhline
    DNS   & Application & 53  & Domain Name System \\ \tblhline
\evenrow
    SMTP  & Application & 25  & Simple Mail Transfer Protocol \\ \tblhline
    TCP   & Transport   & --- & Connection-oriented, reliable \\ \tblhline
\evenrow
    UDP   & Transport   & --- & Connectionless, fast \\ \tblhline
    IP    & Network     & --- & Internet Protocol \\
    \tblbottomrule
    \caption{Common network protocols with layers and port numbers}
    \label{tab:protocols}
\end{longtable}

\begin{longtable}{|L{4cm}|L{4cm}|L{4.5cm}|}
    \longtableheader{3}{Error Detection \& Correction}{%
      \tblcolhdcell{Error Type} & \tblcolhdcell{Detection Method} & \tblcolhdcell{Correction Method}}
    Single-bit error   & Parity check, CRC  & Hamming code \\ \tblhline
\evenrow
    Burst error        & CRC, checksum      & ARQ (retransmit) \\ \tblhline
    Packet loss        & Sequence numbers   & ARQ, FEC \\ \tblhline
\evenrow
    Duplicate frame    & Sequence numbers   & Discard duplicate \\ \tblhline
    Out-of-order frame & Sequence numbers   & Reorder in buffer \\
    \tblbottomrule
    \caption{Error types and handling strategies}
    \label{tab:errors}
\end{longtable}

\begin{longtable}{|C{3cm}|C{3cm}|C{3cm}|C{3cm}|}
    \longtableheader{4}{OSI vs TCP/IP Model}{%
      \tblcolhdcell{OSI Layer} & \tblcolhdcell{TCP/IP Layer} & \tblcolhdcell{PDU Name} & \tblcolhdcell{Key Devices}}
    Application    & Application & Data    & Gateway \\ \tblhline
\evenrow
    Presentation   & Application & Data    & ---     \\ \tblhline
    Session        & Application & Data    & ---     \\ \tblhline
\evenrow
    Transport      & Transport   & Segment & ---     \\ \tblhline
    Network        & Internet    & Packet  & Router  \\ \tblhline
\evenrow
    Data Link      & Link        & Frame   & Switch  \\ \tblhline
    Physical       & Link        & Bit     & Hub     \\
    \tblbottomrule
    \caption{OSI vs TCP/IP model comparison}
    \label{tab:osivsip}
\end{longtable}

% ═══════════════════════════════════════════════════════════════════════════════
\section{Math Shortcuts}
% ═══════════════════════════════════════════════════════════════════════════════

\begin{formulabox}[All Math Commands from \texttt{preamble.tex}]

\textbf{Derivatives:}
\[
    \der{f}{t} \qquad \dder{f}{t} \qquad \pder{f}{x}
\]

\textbf{Vector and matrix notation:}
\[
    \vv{F} = m\vv{a} \qquad \mat{A}\mat{x} = \mat{b}
\]

\textbf{Differential and units:}
\[
    \int_0^T f(t)\dd t \qquad v = 3 \times 10^8\unit{m/s}
\]

\textbf{Magnitude and norm:}
\[
    \abs{\vv{E}} = E_0 \qquad \norm{\mat{A}} = \sqrt{\sum_{ij} a_{ij}^2}
\]

\textbf{Evaluation bar:}
\[
    \eval{\frac{1}{2}mv^2}{0}{T} \qquad \eval{e^{st}}{t=0}{t=\infty}
\]

\textbf{Cancel (strikethrough in math):}
\[
    \frac{\cancel{m} \cdot a}{\cancel{m}} = a
\]
\end{formulabox}

% ═══════════════════════════════════════════════════════════════════════════════
\section{Inline Markers}
% ═══════════════════════════════════════════════════════════════════════════════

\Note This is a note using \texttt{\textbackslash Note}. Use it for side remarks
and clarifications inline in your text.

\medskip
\Imp This is an important callout using \texttt{\textbackslash Imp}. Use it to
flag things the exam will likely test.

\medskip
The command \texttt{\textbackslash term\{...\}} highlights a \term{key term}
being introduced for the first time.

\medskip
\TODO{Replace this paragraph with a proper derivation from the textbook.}

% ═══════════════════════════════════════════════════════════════════════════════
\section{TikZ Block Diagram (Theme-Aware Colors)}
% ═══════════════════════════════════════════════════════════════════════════════

\begin{infobox}[Data Encapsulation Flow]

\begin{center}
\begin{tikzpicture}[
    node distance = 0.8cm,
    block/.style  = {rectangle, draw=defcolor, thick, fill=defcolor!10,
                     minimum width=4.5cm, minimum height=0.7cm,
                     font=\small\bfseries, text=defcolor, align=center},
    arrow/.style  = {->, thick, color=formulacolor}
]
    \node[block] (app)  {Application Layer};
    \node[block] (pres) [below=of app]  {Presentation Layer};
    \node[block] (sess) [below=of pres] {Session Layer};
    \node[block] (trans)[below=of sess] {Transport Layer};
    \node[block] (net)  [below=of trans]{Network Layer};
    \node[block] (dl)   [below=of net]  {Data Link Layer};
    \node[block] (phy)  [below=of dl]   {Physical Layer};

    \draw[arrow] (app)  -- (pres);
    \draw[arrow] (pres) -- (sess);
    \draw[arrow] (sess) -- (trans);
    \draw[arrow] (trans)-- (net);
    \draw[arrow] (net)  -- (dl);
    \draw[arrow] (dl)   -- (phy);

    \node[right=1.2cm of app,  font=\footnotesize\color{darkgray}] {Data};
    \node[right=1.2cm of trans,font=\footnotesize\color{darkgray}] {Segment};
    \node[right=1.2cm of net,  font=\footnotesize\color{darkgray}] {Packet};
    \node[right=1.2cm of dl,   font=\footnotesize\color{darkgray}] {Frame};
    \node[right=1.2cm of phy,  font=\footnotesize\color{darkgray}] {Bits};
\end{tikzpicture}
\end{center}

\end{infobox}

\section{Figure Environment}

\begin{figure}[H]
    \centering
    \safeincludegraphics[width=0.7\textwidth]{example-missing-figure.png}
    \caption{Example figure environment using safe image loading}
    \label{fig:showcase-figure}
\end{figure}

% ═══════════════════════════════════════════════════════════════════════════════
\section{A Real Notes Page — TCP Segment Structure}
% ═══════════════════════════════════════════════════════════════════════════════

This section shows everything working together naturally, the way you'd actually
write notes.

\begin{definitionbox}[TCP — Transmission Control Protocol]
\term{TCP} is a connection-oriented, reliable, byte-stream transport-layer
protocol defined in RFC 793. It guarantees:
\begin{itemize}
    \item Ordered delivery of bytes
    \item Error detection via checksum
    \item Flow control via sliding window
    \item Congestion control (AIMD algorithm)
\end{itemize}
\end{definitionbox}

\subsection{TCP Segment Header Fields}

\begin{longtable}{|L{3.5cm}|C{2cm}|L{7cm}|}
    \longtableheader{3}{TCP Segment Header Fields}{%
      \tblcolhdcell{Field} & \tblcolhdcell{Size (bits)} & \tblcolhdcell{Purpose}}
    Source Port      & 16  & Identifies the sending application \\ \tblhline
\evenrow
    Destination Port & 16  & Identifies the receiving application \\ \tblhline
    Sequence Number  & 32  & Byte offset of first byte in this segment \\ \tblhline
\evenrow
    ACK Number       & 32  & Next expected byte from the other side \\ \tblhline
    Header Length    & 4   & Number of 32-bit words in the header \\ \tblhline
\evenrow
    Control Flags    & 6   & SYN, ACK, FIN, RST, PSH, URG \\ \tblhline
    Window Size      & 16  & Flow control: receiver's buffer space \\ \tblhline
\evenrow
    Checksum         & 16  & Error detection \\ \tblhline
    Urgent Pointer   & 16  & Used when URG flag is set \\ \tblhline
\evenrow
    Options          & var & MSS, timestamps, SACK, window scaling \\
    \tblbottomrule
    \caption{TCP segment header fields}
    \label{tab:tcpheader}
\end{longtable}

\begin{formulabox}[TCP Sequence \& ACK Numbers]
If sender transmits segment starting at byte $n$ with $L$ bytes of data:
\[
    \text{Sequence Number} = n
\]
\[
    \text{Expected ACK from receiver} = n + L
\]
The ACK number always represents the \textbf{next} byte expected, not the last
byte received.
\end{formulabox}

\begin{alertbox}[SYN Flood Attack]
A \term{SYN flood} exploits the TCP 3-way handshake. The attacker sends
many SYN packets with spoofed source IPs. The server allocates resources for
each half-open connection, eventually exhausting memory.

\medskip
\textbf{Mitigation:} SYN cookies — the server does not allocate state until
the handshake completes.
\end{alertbox}

\begin{questionanswerbox}
\textbf{Problem:} Host A sends a TCP segment where Seq\,=\,1000 and the data
is 500 bytes long. What ACK number does Host B send back?

\medskip
\textbf{Solution:}
\[
    \text{ACK} = \text{Seq} + \text{Length} = 1000 + 500 = 1500
\]
Host B's ACK segment will have ACK\,=\,1500, meaning ``I received up to byte
1499; send me byte 1500 next.''
\end{questionanswerbox}

\begin{summarybox}[TCP Key Points]
\begin{itemize}
    \item TCP is \textbf{connection-oriented} — 3-way handshake before data transfer
    \item Reliability via \textbf{sequence numbers + ACKs + retransmission}
    \item Flow control via \textbf{receiver window} (prevents buffer overflow)
    \item Congestion control via \textbf{AIMD} (prevents network collapse)
    \item Connection teardown uses \textbf{4-way FIN handshake}
\end{itemize}
\end{summarybox}

\end{document}
